\documentclass[11pt,a4paper]{article}
\usepackage[utf8]{inputenc}
\usepackage[T1]{fontenc}
\usepackage[spanish]{babel}
\usepackage{geometry}
\usepackage{listings}
\usepackage{hyperref}
\usepackage{booktabs}
\usepackage{float}
\usepackage{amsmath}
\usepackage{amssymb}

\geometry{margin=2.5cm}
\lstset{
    basicstyle=\ttfamily\small,
    breaklines=true,
    frame=single,
    numbers=left,
    numberstyle=\tiny,
    language=C++,
    keywordstyle=\bfseries,
    commentstyle=\itshape
}

\title{Práctico 1 --- Patrones de acceso a memoria}
\author{%
  Pedro Calado Moura \quad Documento: 65092945\\
  Kelvin Martinez \quad Documento: ---\\
  Brandson \quad Documento: ---%
}
\date{}

\begin{document}
\maketitle

%==============================================================================
\section{Plataforma computacional}
%==============================================================================

Todas las librerías de este trabajo se ejecutaron en un sistema operativo Linux con arquitectura de 64 bits en un procesador Intel(R) Core(TM) i7-14700HX (20 núcleos, 28 hilos lógicos) a una frecuencia base de 800 MHz y turbo de hasta 5,5 GHz, y una memoria RAM de 16 GB. Las mediciones se realizaron con \texttt{taskset -c 0} (un único núcleo), así el proceso no salta entre núcleos y los tiempos se interpretan mejor (una sola jerarquía de caché). \textbf{Todos los códigos utilizados para las mediciones fueron compilados sin optimizaciones} (\texttt{g++ -O0} o equivalente).

\textbf{Jerarquía de caché} (salida de \texttt{lscpu -C}, por nivel y una instancia): todas las cachés son \textit{set-associative} (vías \(>1\), conjuntos \(>1\)). L1d (datos): 48\,KiB, 12 vías, 64 conjuntos, línea 64\,B. L1i (instrucciones): 32\,KiB, 8 vías, 64 conjuntos, línea 64\,B. L2 unificada: 2\,MiB, 16 vías, 2048 conjuntos, línea 64\,B. L3 unificada: 33\,MiB, 11 vías, 49\,152 conjuntos, línea 64\,B. Así, cada bloque de memoria se asigna a un conjunto concreto y puede ocupar cualquiera de las vías de ese conjunto.

%==============================================================================
\section{Ejercicio 1 --- Localidad espacial}
%==============================================================================

\subsection{Descripción}

El ejercicio consiste en implementar dos recorridas sobre un arreglo de gran tamaño (100 MB de \texttt{char}) que visitan exactamente la misma cantidad de elementos, pero con patrones de acceso diferentes:

\begin{enumerate}
\item \textbf{Recorrida secuencial:} se visita el arreglo en orden (posición 0, 1, 2, \ldots).
\item \textbf{Recorrida aleatoria:} se visitan las mismas posiciones pero en orden aleatorio.
\end{enumerate}

En ambos casos, el siguiente índice a visitar se obtiene leyendo de un arreglo de índices previamente inicializado. El objetivo es analizar el impacto de la \textbf{localidad espacial} en la jerarquía de memoria: el acceso secuencial explota la pre-carga de bloques contiguos en caché, mientras que el acceso aleatorio tiende a generar múltiples fallos de caché.

\subsection{Resultados}

El programa se ejecutó en un sistema con carga reducida. Ambas recorridas visitan exactamente \(100 \times 1024 \times 1024 = 104\,857\,600\) elementos.

\begin{table}[htbp]
\centering
\begin{tabular}{@{}lc@{}}
\toprule
\textbf{Recorrida} & \textbf{Tiempo (ms)} \\
\midrule
Secuencial         & 260.535              \\
Aleatoria          & 1703.62              \\
\bottomrule
\end{tabular}
\caption{Ejercicio 1: tiempos de ejecución (1 núcleo).}
\label{tab:tiempos}
\end{table}

La recorrida aleatoria resulta aproximadamente \textbf{6.5 veces más lenta} que la secuencial. La variabilidad entre ejecuciones se debe a factores como el estado de la caché al inicio y la carga del sistema.

\subsection{Reflexión}

La diferencia de rendimiento se explica por la \textbf{localidad espacial} y la jerarquía de memoria.

\textbf{Acceso secuencial:} Cuando se accede a posiciones contiguas, la CPU aprovecha las \emph{cache lines} (bloques típicos de 64 bytes). Al leer una posición, se trae a L1/L2 todo el bloque contiguo; los accesos siguientes caen en caché sin necesidad de ir a memoria principal. El ancho de banda efectivo se maximiza.

\textbf{Acceso aleatorio:} Cada acceso puede dirigirse a una dirección distante que no está en caché. Esto provoca un \emph{cache miss} frecuente: la CPU debe ir a RAM para traer el dato. La memoria principal es del orden de 100 veces más lenta que la caché L1, por lo que el tiempo total aumenta drásticamente a pesar de realizar la misma cantidad de operaciones.

En resumen, el experimento demuestra que el \textbf{orden de acceso a los datos} afecta fuertemente el desempeño debido a la jerarquía de memoria. Los patrones de acceso que respetan la localidad espacial (como el secuencial) se benefician del caché y son significativamente más rápidos que los que dispersan los accesos por todo el espacio de direcciones.

%==============================================================================
\section{Ejercicio 2 --- Multiplicación por bloques}
%==============================================================================

\subsection{Ejercicio 2.1 --- Determinación experimental del BS óptimo}

Para tres matrices de tamaño mayor que la capacidad de la caché de último nivel (LLC), se busca el tamaño de bloque \texttt{BS} que maximice el rendimiento (MFLOPS) en la multiplicación por bloques 

Se usó \(N = 1800\), de modo que las tres matrices suman \(3 \times 1800^2 \times 4\,\text{B} \approx 37\,\text{MB} > 33\,\text{MB}\) (LLC). Dado \(N\), conviene restringir \texttt{BS} a divisores de \(N\): si no, aparecen bloques parciales en los bordes, el código debe comprobar \texttt{i < N} en cada iteración interna y el rendimiento empeora. Se hizo un barrido con \texttt{BS\_min=10}, \texttt{BS\_max=200}, \texttt{BS\_step=5} (solo divisores de 1800 en ese rango alineados al step). MFLOPS definidos como \(N^3 / (\text{segundos} \times 10^6)\). Se ejecutó el script cinco veces y se promedió el tiempo para que los resultados sean más representativos; en la tabla se reportan el tiempo medio y los MFLOPS correspondientes.

\begin{table}[H]
\centering
\begin{tabular}{@{}rcc@{}}
\toprule
\textbf{BS} & \textbf{Tiempo (ms)} & \textbf{MFLOPS} \\
\midrule
10  & 13435.15 & 434.88  \\
15  & 13974.39 & 419.80  \\
20  & 13113.64 & 445.58  \\
25  & 12853.51 & 455.57  \\
30  & 13169.35 & 448.87  \\
40  & 12333.37 & 473.44  \\
45  & 12061.91 & \textbf{483.76} \\
50  & 12626.59 & 466.62  \\
60  & 13343.88 & 447.47  \\
75  & 12831.52 & 456.62  \\
90  & 12561.28 & 466.65  \\
100 & 12460.90 & 468.37  \\
120 & 12495.35 & 467.02  \\
150 & 12911.25 & 452.81  \\
180 & 12332.53 & 474.40  \\
200 & 12992.69 & 451.10  \\
\bottomrule
\end{tabular}
\caption{Rendimiento para distintos \texttt{BS} divisores de \(N\) (\(N=1800\), BS\_min=10, BS\_max=200, BS\_step=5, 1 núcleo, \(-O0\)).}
\label{tab:bs}
\end{table}

El \textbf{BS óptimo} en promedio resultó \textbf{45},  Y el BS optmimo de la cinco corridas fueron (45, 60, 90, 60, 60), el promedio de esos mejores BS es 63 y el de los MFLOPS correspondientes 488{,}04. Aunque BS=60 fue el mejor en 3 de 5 corridas, el promedio global favorece a 45 porque en una corrida BS=60 cayó mucho (334{,}71 MFLOPS), lo que baja bastante su media.

\subsection{Comparación con lo esperado teóricamente}

Con la jerarquía de caché ya descrita en la sección~1, se puede estimar un \texttt{BS} teórico. Para que la multiplicación por bloques aproveche la caché, conviene que los tres bloques activos (de \(A\), \(B\) y \(C\)) quepan a la vez. Con \texttt{float} (4\,B por elemento), el tamaño total es \(3 \cdot BS^2 \cdot 4 = 12\,BS^2\) bytes. Usando como referencia la L1d por núcleo (48\,KiB según \texttt{lscpu -C}):

\[
12\,BS^2 \leq 48\,\text{KiB} \quad \Rightarrow \quad BS \lesssim \sqrt{\frac{48 \cdot 1024}{12}} \approx 64.
\]

Posibles desvíos entre teoría y práctica se explican porque el modelo es idealizado: no toda la caché (L1) se dedica solo a los tres bloques, esta ocupada por otras líneas de caché activas , pila y otras estructuras; y  el acceso a \(B[kN+j]\) por columnas no siempre tiene localidad óptima; y además influyen la asociatividad, la política de reemplazo, el prefetching y el coste extra de los bucles de bloqueo. En conjunto, los resultados experimentales no coinciden exactamente con la estimación teórica: el mejor promedio se obtuvo con \texttt{BS = 45}. Sin embargo, \texttt{BS = 60} también mostró un desempeño muy competitivo y fue el mejor en varias ejecuciones, por lo que los resultados siguen siendo compatibles, de manera aproximada, con la predicción basada en la jerarquía de caché

\subsection{Ejercicio 2.2 --- Variante lineal ikj y orden de los bucles por bloques}

\subsection{Comparación entre la versión lineal y la versión por bloques}

Al inspeccionar el código de la multiplicación lineal, se observa que el orden de loops \texttt{ikj} es más conveniente que \texttt{ijk} para matrices almacenadas en \emph{row-major}. La razón principal es que, con \texttt{ikj}, el índice \texttt{j} queda en el loop más interno, por lo que tanto \texttt{B[k*N + j]} como \texttt{C[i*N + j]} se recorren de forma contigua en memoria. Además, el valor \texttt{A[i*N + k]} se carga una sola vez y se reutiliza para todas las posiciones de \texttt{j}. En cambio, con el orden \texttt{ijk}, el índice \texttt{k} queda en el loop interno, de modo que el acceso a \texttt{B[k*N + j]} avanza por columnas de \texttt{B}, lo cual implica un \emph{stride} de tamaño \(N\) y, por tanto, una peor localidad espacial.

A partir de esta observación, nuestro grupo intentó trasladar la misma idea a la versión bloqueada. En lugar del orden original \texttt{ii,jj,kk,i,j,k}, se probó reordenar los loops como \texttt{ii,kk,jj,i,k,j}, buscando que dentro de cada bloque el recorrido mantuviera a \texttt{j} como índice interno, igual que en la versión lineal \texttt{ikj}.

La Tabla~\ref{tab:comparacion_tres} muestra el promedio de cinco ejecuciones para la versión bloqueada original, la versión bloqueada con loops reordenados y la versión lineal. En todos los casos se utilizó \(N=1800\), con \(BS=45\), fijando la ejecución a un único núcleo mediante \texttt{taskset -c 0}.

\begin{table}[ht]
\centering
\caption{Promedio de cinco ejecuciones: comparación entre la versión bloqueada original, la bloqueada reordenada y la lineal.}
\label{tab:comparacion_tres}
\begin{tabular}{|l|r|r|}
\hline
Versión & Tiempo promedio (ms) & MFLOPS promedio \\ \hline
Bloqueado \texttt{(ii,jj,kk)} & 12882.33 & 453.89 \\ \hline
Bloqueado \texttt{(ii,kk,jj)} & 11851.43 & 492.19 \\ \hline
Lineal \texttt{ikj} & 10384.74 & 561.72 \\ \hline
\end{tabular}
\end{table}

Los resultados muestran que el simple reordenamiento de loops en la variante bloqueada sí mejora el desempeño respecto de la versión bloqueada original. En promedio, \texttt{(ii,kk,jj)} reduce el tiempo de ejecución y aumenta el rendimiento en MFLOPS frente a \texttt{(ii,jj,kk)}. Sin embargo, aun con esa mejora, la variante lineal \texttt{ikj} sigue siendo superior.

\subsection{Ejercicio 2.3 --- Comparación de órdenes de loop lineales}

Se comparó el tiempo de ejecución y el rendimiento en MFLOPS (\(N^3 / (\text{segundos} \times 10^6)\)) de la versión \emph{lineal} (sin bloques) del producto de matrices para los seis órdenes de loop posibles: \texttt{ijk}, \texttt{jik}, \texttt{ikj}, \texttt{kij}, \texttt{jki} y \texttt{kji}. Se utilizaron los tamaños de matriz 32, 128, 256, 260, 512, 550, 1024 y 1050, ejecutando cada combinación cuatro veces y promediando los MFLOPS para reducir la variabilidad. La Tabla~\ref{tab:loops} resume los resultados.

\begin{table}[htbp]
\centering
\begin{tabular}{@{}rcccccc|c@{}}
\toprule
\textbf{N} & \textbf{ijk} & \textbf{jik} & \textbf{ikj} & \textbf{kij} & \textbf{jki} & \textbf{kji} & \textbf{Prom. N} \\
\midrule
32   & 510.25 & 504.75 & \textbf{603.00} & 516.00 & 506.25 & 500.00 & 522.21 \\
128  & 523.00 & 517.75 & \textbf{669.00} & 566.00 & 501.25 & 498.25 & 545.88 \\
256  & 525.00 & 510.00 & \textbf{658.25} & 531.25 & 449.00 & 435.25 & 518.29 \\
260  & 483.25 & 482.75 & \textbf{630.25} & 525.00 & 494.25 & 488.25 & 517.29 \\
512  & 471.00 & 470.75 & \textbf{633.75} & 547.00 & 414.00 & 418.50 & 492.38 \\
550  & 486.75 & 490.75 & \textbf{584.25} & 502.00 & 413.00 & 397.50 & 479.04 \\
1024 & 301.00 & 282.50 & \textbf{633.50} & 566.25 & 193.50 & 181.00 & 359.63 \\
1050 & 410.50 & 450.75 & \textbf{642.75} & 520.00 & 321.50 & 336.75 & 447.04 \\
\midrule
Promedio total & 408.78 & 461.28 & \textbf{631.09} & 534.19 & 411.59 & 408.19 & 482.72 \\
\bottomrule
\end{tabular}
\caption{MFLOPS promedio (4 ejecuciones) de la versión lineal para distintos órdenes de loop (1 núcleo, \(-O0\)).}
\label{tab:loops}
\end{table}

\begingroup
\setlength{\parindent}{0pt}
\setlength{\parskip}{0.7em}
\sloppy

Los resultados obtenidos muestran una tendencia clara: las variantes \texttt{ikj} y \texttt{kij} son, en general, las de mejor desempeño, mientras que \texttt{jki} y \texttt{kji} resultan las peores, especialmente para tamaños grandes. Esto puede explicarse a partir del patrón de acceso a memoria y de la jerarquía de caché de la máquina utilizada.

En esta arquitectura, los accesos contiguos por fila aprovechan mejor la localidad espacial: una línea de caché trae varios elementos consecutivos, que luego pueden reutilizarse antes de ser reemplazados. Por eso, cuando el índice \texttt{j} queda en el loop más interno, como en \texttt{ikj} y \texttt{kij}, se recorren de forma contigua las filas de \(B\) y \(C\), lo que permite un uso más eficiente de las líneas de caché.

La variante \texttt{ikj} resulta, en promedio, la mejor. Como las matrices se almacenan por filas (como se comentó en clase), en este orden, para cada par \((i,k)\), el valor \(A[i,k]\) se lee una sola vez y se reutiliza durante todo el recorrido de \texttt{j}. Además, tanto \(B[k,j]\) como \(C[i,j]\) se recorren secuencialmente en memoria. Esto combina buena localidad espacial con reutilización temporal, y explica por qué su rendimiento se mantiene alto incluso cuando el tamaño del problema supera L1 y L2. La variante \texttt{kij} también deja a \texttt{j} en el loop interno, por lo que conserva un acceso contiguo favorable, aunque con menor reutilización local sobre \(C\) que \texttt{ikj}.

En las variantes \texttt{ijk}, \texttt{jik}, \texttt{jki} y \texttt{kji}, el problema principal es que el loop interno obliga a recorrer por columnas, lo que empeora mucho la localidad espacial. En \texttt{ijk} y \texttt{jik}, esto afecta sobre todo a \(B[k,j]\), ya que al variar \texttt{k} se avanza por columnas de \(B\), produciendo saltos de tamaño \(N\). En \texttt{jki} y \texttt{kji}, la situación es todavía peor, porque al variar \texttt{i} se recorren por columnas tanto \(A[i,k]\) como \(C[i,j]\). En todos estos casos, el procesador trae líneas completas de caché de 64\,B, pero solo se aprovecha una pequeña parte antes de saltar a otra región lejana de memoria. Esto genera más misses de caché y una muy baja reutilización efectiva.

También se observa que para tamaños pequeños, como \(N=32\), las diferencias entre órdenes existen pero son moderadas. Esto es razonable, ya que el conjunto de trabajo es pequeño y entra en L1. En cambio, para \(N=128\) el conjunto de trabajo ya supera L1 pero todavía entra en L2, y para \(N=256\), \(260\), \(512\), \(550\), \(1024\) y \(1050\) supera L2 pero entra en L3. A medida que el tamaño de las matrices deja de caber en los niveles más rápidos de caché, el costo de un mal patrón de acceso se vuelve mucho más visible, y por eso las diferencias entre órdenes crecen marcadamente.

Un detalle interesante es la diferencia entre \(N=1024\) y \(N=1050\).

\endgroup
\end{document}